% ======================================================================
%  sec21.tex —— 第 21 节 不可约概形
%  底本：原书 pp.105–107（印刷页 = PDF 页）
%  原书标题：21. Irreducible schemes
%            21.1 Irreducible spaces / 21.2 Irreducible components
%            21.3 Irreducible schemes / 21.4 Integral schemes
% ======================================================================

\section{不可约概形}\label{sec:21}

% ---------------------------------------------------------------- p.105
\subsection{不可约空间}\label{sec:21.1}

\begin{definition}\label{def:21.1.1}
称拓扑空间 $X$ \emph{不可约}，若它非空且不是两个真闭子空间之并。
\end{definition}

\begin{yizhu}
Stacks 把同一件事写成：$X$ 非空，且只要 $X=Z_1\cup Z_2$ 而 $Z_1,Z_2$ 闭，就有 $X=Z_1$ 或 $X=Z_2$。两种说法逻辑等价，但后者便于对有限多个闭子空间归纳：不可约空间不能写成有限个真闭子空间之并。定义中「非空」一项不可省，否则空空间按上述条件会被判为不可约，而空概形是既约的，这将使「整概形」的定义失去意义。
[来源：Stacks Project, Tag 004V（Definition 5.8.1, Irreducible components），\url{https://stacks.math.columbia.edu/tag/004V}（2026-11-17 访问）。]
\end{yizhu}

\begin{exercise}\label{exc:21.1.2}
设 $X$ 是不可约空间。则
\begin{enumerate}[label=(\roman*),leftmargin=2em,itemsep=0.25ex]
  \item 空间 $X$ 连通。
  \item $X$ 的任意非空开子集不可约。
  \item $X$ 的任意非空开子集在 $X$ 中稠密。
  \item $X$ 的两个非空开子集之交非空。
\end{enumerate}
\end{exercise}

\begin{exercise}\label{exc:21.1.3}
拓扑空间中不可约子空间的闭包不可约。
\end{exercise}

\begin{yizhu}
Stacks 的证明只有两行：设 $T\subseteq X$ 不可约，若其闭包 $\overline{T}=Z_1\cup Z_2$ 是闭集之并，则 $T=(T\cap Z_1)\cup(T\cap Z_2)$，而各 $T\cap Z_i$ 在 $T$ 中闭，故 $T$ 等于其中之一，再取闭包即得 $\overline{T}$ 等于对应的 $Z_i$。可见关键是「在 $T$ 中闭」而非「在 $X$ 中闭」，这一步在仿射概形的语言里对应「取零点集后仍可回到素理想」。
[来源：Stacks Project, Tag 004W（Lemma 5.8.3, Irreducible components），\url{https://stacks.math.columbia.edu/tag/004W}（2026-11-17 访问）。]
\end{yizhu}

\begin{exercise}\label{exc:21.1.4}
不可约空间在连续映射下的像不可约。
\end{exercise}

\begin{definition}\label{def:21.1.5}
称拓扑空间 $X$ 有\emph{一般点} $x\in X$，若 $X=\overline{\{x\}}$。
\end{definition}

\begin{yizhu}
Stacks 只对不可约闭子集 $Z\subseteq X$ 定义一般点：点 $\xi\in Z$ 称为 $Z$ 的一般点，若 $Z=\overline{\{\xi\}}$。原书对空间本身定义（$X=\overline{\{x\}}$），正是 $Z=X$ 的特例。两种措辞的差别不只是编排：对子集定义时「$\xi$ 落在 $Z$ 中」是条件的一部分，这解释了为何一般点唯一时它必属于它所代表的闭集。
[来源：Stacks Project, Tag 004X（Definition 5.8.6, Irreducible components），\url{https://stacks.math.columbia.edu/tag/004X}（2026-11-17 访问）。]
\end{yizhu}

\begin{exercise}\label{exc:21.1.6}
若空间 $X$ 有一般点，则 $X$ 不可约。
\end{exercise}

\begin{exercise}\label{exc:21.1.7}
若 $X$ 是 Kolmogorov 空间，则 $X$ 的任意两个一般点相等。
\end{exercise}

\subsection{不可约分支}\label{sec:21.2}

\begin{definition}\label{def:21.2.1}
设 $X$ 是拓扑空间。$X$ 的\emph{不可约分支}是指 $X$ 的极大不可约子空间。
\end{definition}

\begin{exercise}\label{exc:21.2.2}
$X$ 的任意不可约分支闭。
\end{exercise}

\begin{exercise}\label{exc:21.2.3}
$X$ 的全部不可约分支之并等于 $X$。
\end{exercise}

\subsection{不可约概形}\label{sec:21.3}

\begin{definition}\label{def:21.3.1}
称概形 $X$ \emph{不可约}，若其底拓扑空间不可约。
\end{definition}

\begin{lemma}\label{lem:21.3.2}
设 $X$ 是概形。则 $X$ 的任意闭不可约子空间有唯一的一般点。
\end{lemma}

\begin{yizhu}
唯一性只用 Kolmogorov 性质（习题 \ref{exc:21.1.7}）：若 $x\ne y$ 而闭包相同，则任何开集含 $x$ 当且仅当含 $y$，与 $T_0$ 矛盾。存在性在仿射情形化为一句交换代数：$\Spec(A)$ 中素理想的闭包是 $V(\mathfrak{p})$，而每个闭不可约子集都形如 $V(\mathfrak{p})$。书中绕开这一步，改用「换 $X$ 为 $\Spec(A/I)$ 并取既约」的手法（见下方证明），于是只需证既约不可约环的零理想为素。
[来源：Stacks Project, Tag 00ES（Lemma 10.26.1, Irreducible components of spectra），\url{https://stacks.math.columbia.edu/tag/00ES}（2026-11-17 访问）。]
\end{yizhu}

\begin{proof}
唯一性由习题 \ref{exc:21.1.7} 得到。

下证存在性。设 $Y\subseteq X$ 是闭不可约子空间。取仿射开子集 $U\subseteq X$ 使 $Y\cap U$ 非空。因 $Y$ 不可约，有 $Y=\overline{Y\cap U}$。于是只需证 $Y\cap U$ 有一般点。故可把 $X$ 换成 $U$，从而设 $X=\Spec(A)$ 仿射。

取根理想 $I\subseteq A$ 使 $Y=Z(I)$。只需证 $I$ 是素理想，因为此时对应的点将给出 $Z(I)$ 的一般点。把 $X$ 换成 $\Spec(A/I)$，可设 $I=(0)$，从而 $\Spec(A)$ 既约且不可约。

现在设 $xy\in(0)$。注意 $\Spec(A)$ 被 $Z(x)$ 与 $Z(y)$ 覆盖。因 $\Spec(A)$ 不可约，不妨设 $Z(x)=\Spec(A)$。这蕴含 $\sqrt{(x)}=\sqrt{(0)}$。

% ---------------------------------------------------------------- p.106
因 $\Spec(A)$ 既约，必有 $x=0$。这说明 $(0)$ 是素理想，如所愿。
\end{proof}

\begin{corollary}\label{cor:21.3.3}
设 $X$ 是概形。则 $x\mapsto\overline{\{x\}}$ 给出 $X$ 的底层集与 $X$ 的闭不可约子空间集之间的双射。
\end{corollary}

\begin{corollary}\label{cor:21.3.4}
$\Spec(A)$ 的不可约分支集与 $A$ 的极小素理想集一一对应。
\end{corollary}

\begin{corollary}\label{cor:21.3.5}
仿射概形 $\Spec(A)$ 不可约当且仅当 $A$ 的幂零根是素理想。
\end{corollary}

\begin{exercise}\label{exc:21.3.6}
设 $X$ 是概形，$x\in X$ 是点。下列条件等价：
\begin{itemize}[leftmargin=1.5em,itemsep=0.15ex]
  \item 存在唯一的含 $x$ 的 $X$ 的不可约分支。
  \item $\cO_{X,x}$ 的幂零根是素理想。
\end{itemize}
\end{exercise}

\begin{lemma}\label{lem:21.3.7}
设 $X$ 是概形。下列条件等价：
\begin{enumerate}[label=(\roman*),leftmargin=2em,itemsep=0.25ex]
  \item 概形 $X$ 不可约。
  \item 概形 $X$ 非空，且 $X$ 的任意非空仿射开子概形 $U$ 不可约。
  \item 存在非空的覆盖 $X=\bigcup_{i\in I}U_i$，使每个 $U_i$ 是不可约仿射开子概形，且每个交 $U_i\cap U_j$ 非空。
\end{enumerate}
\end{lemma}

\begin{proof}
(1)$\Rightarrow$(2) 由习题 \ref{exc:21.1.2}(2) 得到。

(2)$\Rightarrow$(3)：取任意覆盖 $X=\bigcup_{i\in I}U_i$，使每个 $U_i$ 是非空仿射开子概形。按假设，每个 $U_i$ 不可约。若某对 $i,j$ 使 $U_i\cap U_j$ 为空，则不交并 $U_i\sqcup U_j$ 也是 $X$ 的非空仿射开子概形，却不可约，因为 $U_i$ 与 $U_j$ 都是 $U_i\sqcup U_j$ 的闭子空间，这与 (2) 矛盾。故 $U_i\cap U_j$ 非空，如所愿。

(3)$\Rightarrow$(1)：设底拓扑空间 $X$ 是两个闭子空间之并 $Z_1\cup Z_2$。只需证 $X=Z_1$ 或 $X=Z_2$。因每个 $U_i$ 不可约，必有 $U_i\subseteq Z_1$ 或 $U_i\subseteq Z_2$。不妨取 $i\in I$ 使 $U_i\subseteq Z_1$。对每个 $j\in I$，注意 $U_i\cap U_j$ 在 $U_j$ 中稠密（习题 \ref{exc:21.1.2}(3)），且含于闭子集 $Z_1\cap U_j$。于是必有 $U_j=Z_1\cap U_j$，即 $U_j\subseteq Z_1$。这说明 $X=Z_1$，如所愿。
\end{proof}

% ---------------------------------------------------------------- p.106
\subsection{整概形}\label{sec:21.4}

\begin{definition}\label{def:21.4.1}
称概形 $X$ \emph{整}，若它不可约且既约。
\end{definition}

\begin{yizhu}
Stacks 的定义取另一条等价路径：$X$ 整，指它非空且每个非空仿射开子概形 $\Spec(R)$ 的 $R$ 是整环。「不可约 + 既约」与这条要经过引理 \ref{lem:21.3.7} 与推论 \ref{cor:19.1.3} 才能对上，其中不可约供出素理想、既约供出零理想，两者合起来才使 $R$ 成为整环。书中的定义更短，但把「非空」藏进了不可约。
[来源：Stacks Project, Tag 01OK（Definition 28.3.1, Integral, irreducible, and reduced schemes），\url{https://stacks.math.columbia.edu/tag/01OK}（2026-11-17 访问）。]
\end{yizhu}

结合推论 \ref{cor:21.3.5} 与推论 \ref{cor:19.1.3}，得到下述结果：

\begin{corollary}\label{cor:21.4.2}
仿射概形 $\Spec(A)$ 整当且仅当 $A$ 是整环。
\end{corollary}

\begin{exercise}\label{exc:21.4.3}
设 $X$ 是概形。下列条件等价：
\begin{enumerate}[label=(\roman*),leftmargin=2em,itemsep=0.25ex]
  \item 概形 $X$ 整。
  \item 概形 $X$ 非空，且对任意开子概形 $U\subseteq X$，环 $\cO_X(U)$ 是整环。
% ---------------------------------------------------------------- p.107
  \item 概形 $X$ 非空，且对任意仿射开子概形 $U\subseteq X$，环 $\cO_X(U)$ 是整环。
\end{enumerate}
\end{exercise}
